Fix \(z\) and let \(q(\cdot,z)\) be a common conditional density with \(c\le q\le C\).  For a direction \(h(\cdot,z)\) satisfying \(\int_0^1h(y,z)dy=0\), define
\[
H_h(y,z)=\int_0^y h(x,z)dx .
\]
Then \(H_h(0,z)=H_h(1,z)=0\).  Consider the symmetric path \(q_{0,t}=q-th/2\), \(q_{1,t}=q+th/2\), for \(|t|\) small enough that both densities stay positive.  If \(F_{a,t}\) and \(Q_{a,t}=F_{a,t}^{-1}\) are the corresponding distribution and quantile functions, differentiation of \(F_{a,t}\{Q_{a,t}(u,z),z\}=u\), justified by \(q_{a,t}\ge c/2\), gives
\[
\dot Q_{1,0}(u,z)-\dot Q_{0,0}(u,z)
=-\frac{H_h(Q(u,z),z)}{q(Q(u,z),z)}.                                      \tag{1}
\]
At \(t=0\) the two quantiles coincide.  Hence the first derivative of their squared \(L^2(0,1)\) distance is zero, and differentiating a second time gives
\[
\frac12D^2W_2^2(q\,dy,q\,dy)[h,h]
=\int_0^1\{\dot Q_{1,0}(u,z)-\dot Q_{0,0}(u,z)\}^2du.
\]
Changing variables \(u=F(y,z)\) in this display and using (1) yields
\[
\frac12D^2W_2^2(q\,dy,q\,dy)[h,h]
=\int_0^1\frac{H_h(y,z)^2}{q(y,z)}dy.                                      \tag{2}
\]
Let \(u_h=\mathcal G_{n,z}h\).  Integrating the defining equation
\(-\partial_y(q\partial_yu_h)=h\) from \(0\) to \(y\), and using the Neumann boundary condition at zero, gives
\[
q(y,z)\partial_yu_h(y,z)=-H_h(y,z).
\]
Integration by parts has no boundary term because \(H_h(0,z)=H_h(1,z)=0\), and therefore
\[
\int_0^1h(y,z)u_h(y,z)dy
=\int_0^1q(y,z)|\partial_yu_h(y,z)|^2dy
=\int_0^1\frac{H_h(y,z)^2}{q(y,z)}dy.                                      \tag{3}
\]
Equations (2)--(3), integrated with respect to the common marginal \(R_n\), prove the asserted vanishing first derivative and inverse-elliptic second variation.  The common marginal is supplied by \(\mathrm{ass{:}shared\text{-}covariate\text{-}marginal}\), and every division above is licensed by \(\mathrm{ass{:}conditional\text{-}density\text{-}bounds}\).

For the displayed cosine pair, put \(e=b(z)/2\).  The two conditional distribution functions are
\[
F_{1,e}(y)=y+\frac{e}{2\pi}\sin(2\pi y),\qquad
F_{0,e}(y)=y-\frac{e}{2\pi}\sin(2\pi y).
\]
For \(|e|<1\), their derivatives are uniformly bounded away from zero.  Taylor expansion of the inverse identities, with the remainder controlled uniformly by that lower derivative bound, gives
\[
F_{1,e}^{-1}(u)-F_{0,e}^{-1}(u)
=-\frac{2e}{2\pi}\sin(2\pi u)+O(|e|^2)
=-\frac{b(z)}{2\pi}\sin(2\pi u)+O(|b(z)|^2).                              \tag{4}
\]
Squaring (4), integrating over \(u\), and using \(\int_0^1\sin^2(2\pi u)du=1/2\) proves
\[
W_2^2(Q_{0,n}(\cdot\mid z),Q_{1,n}(\cdot\mid z))
=\frac{b(z)^2}{8\pi^2}+O(|b(z)|^3).                                        \tag{5}
\]
Since \(R_n\) is a probability measure, its integral remainder is \(O(\|b\|_\infty^3)\).  Finally, on the one-sided path \(q_0=1\), \(q_1=1+t b(z)\cos(2\pi y)\), equation (2) applies with \(H_h=b(z)\sin(2\pi y)/(2\pi)\).  Thus the coefficient of \(t^2\) is \(\int b^2dR_n/(8\pi^2)\), and the second derivative is \(\int b^2dR_n/(4\pi^2)\), as claimed.